\documentclass[10pt,a4paper]{amsart}
\usepackage[margin=19mm]{geometry}
\usepackage{amsmath,amssymb,mathtools}
\usepackage[scr=rsfs,cal=euler]{mathalfa}
\usepackage{xfrac}
\usepackage[unicode,hidelinks,bookmarks=false]{hyperref}
\newcommand{\ie}{i.e.\ }
\newcommand{\cf}{cf.\ }
\newcommand{\resp}{respectively\ }
\newcommand{\Subsection}{\subsection}
\providecommand{\nobreakdash}{\nobreak\hbox{-}}
\setlength{\emergencystretch}{2em}
\multlinegap0pt
\usepackage{bm}
\usepackage{bbm}
\usepackage{dsfont}
\usepackage{xspace}
\usepackage{cancel}
\usepackage{mathtools}
\usepackage{amsthm}
\makeatletter
\@ifundefined{thm}{\newtheorem{thm}{Thm}}{}
\@ifundefined{lem}{\newtheorem{lem}[thm]{Lemma}}{}
\@ifundefined{prop}{\newtheorem{prop}[thm]{Proposition}}{}
\makeatother
\def\ch{{\mathcal H}}
\def\ga{{\mathfrak A}} \def\gpa{{\mathfrak a}}
\def\om{\omega} \def\Om{\Omega}
\title[Quasi-invariant states for the action of compact groups]{Quasi-invariant states for the action of compact groups}
\author{Maria Elena Griseta}
\date{Source: arXiv:2507.20931v1 (CC BY 4.0)}
\begin{document}
\maketitle

\noindent\emph{Source and licence.} Quasi-invariant states for the action of compact groups. Version arXiv:2507.20931v1 (CC BY 4.0), distributed under the Creative Commons Attribution 4.0 International licence, as recorded in the arXiv licence metadata for this exact version and shown on the abstract page https://arxiv.org/abs/2507.20931v1. Full rights notice: licenses/arxiv-2507.20931-CC-BY-4.0.txt.\par\smallskip

\noindent\textit{Editorial scope.} Counted unit: the Prop whose source label is \texttt{prop:isom} in the article (source0.tex lines 425--428, source label \texttt{prop:isom}), delivered together with the article's own complete original proof of it (source0.tex lines 430--449, one \texttt{proof} environment of 20 lines). No mathematics is written, repaired or re-derived: statement and proof are the article's own text line for line. Required context retained uncounted: lem:affiliation (lines 393--399); lem:nonsingular (lines 413--416). Every definition, display and label of those lines is retained unchanged. Omitted from this package and not redistributed: 1--392, 400--412, 417--424, 429--429, 450--502. Nothing inside a retained excerpt is omitted or altered: every retained body is delivered exactly as the article's lines are written, so no omission receipt of any kind accompanies this package. Citations: the retained excerpts carry no citation command at all, so no bibliography entry of the article is transcribed and no reference is redistributed. Reference adaptation: none was needed and none was made; every reference inside the delivered unit resolves inside it, so no unresolved reference or citation is produced. Printed slips kept literal: no printed word, symbol, hypothesis or number of the counted statement or of its proof was altered. Layout and class: the article's own class is \texttt{amsart} with packages including amsmath, amsthm, amsfonts, eucal, xypic, graphicx, amssymb, xcolor, url; the wrapper is the lane's pinned \texttt{amsart} editorial wrapper at 10pt on A4, which restates the theorem-like environments the retained text needs and adds the article's own macro definitions that the retained text needs (\texttt{\textbackslash ch}, \texttt{\textbackslash ga}, \texttt{\textbackslash om}), with extra wrapper packages (bm, bbm, dsfont, xspace, cancel, mathtools). The delivered numbering is the unit's own. Assets: no retained span contains a figure environment, a graphics-inclusion command, a PDF-page inclusion, a table or a TikZ picture; no image file is copied into this package, the package declares no asset dependency and no image file is redistributed with it. Licence: version arXiv:2507.20931v1 of this article is distributed under the Creative Commons Attribution 4.0 International licence, as recorded in the arXiv licence metadata for this exact version and shown on the abstract page https://arxiv.org/abs/2507.20931v1. A full rights notice is delivered at licenses/arxiv-2507.20931-CC-BY-4.0.txt. Characters: every retained excerpt is pure ASCII, so the delivered engine, XeLaTeX with its default Latin Modern text font, reports no missing character.
\begin{lem}
\label{lem:affiliation}
If $T: D(T)\subset \ch\rightarrow\ch$ is a non-singular self-adjoint operator, then its inverse
$T^{-1}$, defined on its natural domain $D(T^{-1})={\rm Ran}(T)$ is a self-adjoint operator.
Moreover, if $T$ is affiliated to a von Neumann algebra $\mathcal{R}$, then $T^{-1}$ is also
affiliated to $\mathcal{R}$.
\end{lem}

\begin{lem}
\label{lem:nonsingular}
The operator $H$ defined as the Radon-Nikodym derivative of $\varphi_\om$ w.r.t $\om$ is non-singular.
\end{lem}

\begin{prop}
\label{prop:isom}
The isometry $W$ is surjective and thus unitary. 
\end{prop}

\begin{proof}
Since the range of $W$ contains the subspace $\{\pi_\om(a) H^\frac{1}{2}\xi_\om: a\in\ga\}$, it is enough
to prove that $H^\frac{1}{2}\xi_\om$ is cyclic for $\pi_\om(\ga)$. By von Neumann's bicommutant theorem, this is the same as 
showing that  $H^\frac{ 1}{2}\xi_\om$ is cyclic for $\mathcal{R}_\om= \pi_\om(\ga)''$. Set $T=H^\frac{1}{2}$ and note 
that $T$ is non-singular by virtue of Lemma \ref{lem:nonsingular}.\\
Let now $x$ be any vector in $\ch_\om$. For every $\varepsilon >0$, by cyclicity of
$\xi_\om$ there exists
$A$ in $\mathcal{R}_\om$ such that $\|A\xi_\om- x \|<\frac{\varepsilon}{2}$.
Thanks to Lemma \ref{lem:affiliation}, $T^{-1}$ is a (positive) self-adjoint operator affiliated with $\mathcal{R}_\om$ as it is in fact affiliated with the centraliser of $\om$. For every $n>0$, let $E_n$ be the spectral
projection of $T^{-1}$ corresponding to the interval $[0, n]$. Clearly, $E_n$ still sits in 
$\mathcal{R}_\om$, which means $B_n:=AE_nT^{-1}$ is a bounded operator that belongs to $\mathcal{R}_\om$.
Because $E_n$ strongly converges to the identity $I$ as $n$ goes to infinity, there must exist $n_0$ such that
$\|E_{n_0}\xi_\om- \xi_\om\|<\frac{\varepsilon}{2\|A\|}$.
But then we have
\begin{align*}
\|B_{n_0}T\xi_\om- x\|\leq&\,\|AE_{n_0}T^{-1}T\xi_\om- A\xi_\om\|+\|A\xi_\om-x\| \leq\\
 &\,\|A\|\|E_{n_0}\xi_\om- \xi_\om\|+\|A\xi_\om-x\| \leq \frac{\varepsilon}{2}+\frac{\varepsilon}{2}=\varepsilon\,,
\end{align*}
which ends the proof as  $\varepsilon$ is arbitrary.
\end{proof}

\end{document}
